\section{Conclusion}\label{sec:conclusion}

This paper was motivated by a critical open question in scientific machine learning: are the catastrophic failures of Physics-Informed Neural Networks (PINNs) on stiff partial differential equations the result of suboptimal hyperparameter tuning, or are they embedded within the fundamental architecture itself? To rigorously answer this, we constructed an experimental Atlas comprising 24 controlled experiments spanning three canonical PDEs—the 1D Advection equation, Viscous Burgers, and the 2D Helmholtz problem—evaluated against an array of robust optimization and physical diagnostic metrics. The ensemble of results presented in this work demonstrates that PINN failure under high stiffness is fundamentally a structural pathology rather than a superficial artifact of empirical tuning.

Our investigation isolated the spectral barrier as a primary structural obstacle. By synthesizing the results of Exps.~1–3 and the phase mapping in Exp.~25, we established that the stiffness-induced spectral bias is quantitatively measurable at initialization via the Neural Tangent Kernel conditioning. Furthermore, this barrier is persistently present across all standard activation functions, strictly independent of internal network capacity scaling, and manifests as a severe performance cliff in the 2D architectural phase space—the network transitions instantaneously from trivially solvable to structurally impossible, with no intermediate scaling regime where adding collocation points or network width continuously rescues the solver. Taken together, these results indicate that the spectral failure is embedded in the representational geometry of the MLP architecture, not in any specific optimization choice.

Beyond representation limits, we identified profound pathologies in the associated optimization landscapes. Through Exps.~5, 13–15, and 27, we observed that as PDE stiffness escalates, the loss landscape actively directs the gradient-based optimizer toward geometrically sharp, physically spurious traps that effectively bypass the target physics. While intensive hyperparameter tuning provides local navigation within this hostile landscape, it systematically fails to restructure it. The variance decomposition analysis corroborates this phenomenon: for the Viscous Burgers equation, the structural factor group accounts for an estimated 292 times more variance (upper bound given the experimental design; see Section~9) compared to standard hyperparameter interventions, decisively cementing the structural dominance hypothesis for diffusion-dominated failure. This structural dominance is strongest for diffusion-dominated regimes (Burgers: $\etasq{}(S)/\etasq{}(H) \approx 292$) and more moderate for purely hyperbolic cases (Advection: $\etasq{}(S)/\etasq{}(H) \approx 1.7$), suggesting that the degree of structural determination scales with the dissipative character of the PDE operator.

Crucially, our analysis establishes that the dominant failure mechanisms are largely independent in their structural origin, despite their consistent temporal co-occurrence. While the instantaneous correlation between spectral and gradient failure signals is moderate ($r=0.348$, Exp.~18), the lagged cross-correlation peaks at $r=0.630$ with a temporal offset of $\sim 600$ training steps, indicating that gradient pathology's diagnostic signature precedes the deepening of spectral error. This sequential onset pattern is consistent with both mechanisms being independently triggered by the same underlying stiffness, rather than one directly causing the other. The mode-specific recovery profiles of Experiment~22.2 provide independent corroboration: interventions targeting gradient pathology do not incidentally resolve spectral bias, and vice versa, confirming structural separability. The ``Zugzwang'' analogy describes not a forced coupling between modes, but rather the absence of any escape move: both modes are simultaneously active, and the available interventions address neither structurally. Correcting the spectral bias alone does not resolve the gradient pathology, and mitigating gradient conflict does not bypass the spectral barrier.

Even when PINNs successfully minimize the localized PDE residuals and achieve low scalar $\Ltwo$ errors, they frequently remain physically discontinuous. The conservation law audit in Exp.~26 illustrates that $\Ltwo$ accuracy is mathematically necessary but fundamentally insufficient for physical fidelity, as seemingly successful models can exhibit catastrophic global mass and momentum drift. For practical engineering deployments, this dictates that PINN solutions must be continuously validated against macroscopic global integral invariants, rather than relying solely on pointwise optimization convergence.

Addressing these foundational vulnerabilities requires a specific, mechanistically grounded path forward that abandons the paradigm of unconstrained post-hoc tuning. Hard-constrained architectures that satisfy boundary and initial conditions exactly by construction \cite{Lagaris1998, Berg2018} offer a direct mathematical route to bypassing the temporal boundary-value conflicts responsible for trivial collapse. Furthermore, transitioning from scalar function approximators to operator-based representations, such as DeepONets \cite{Lu2021deeponet} or Fourier Neural Operators \cite{Li2021fno}, may restructure the function space sufficiently to alleviate the NTK spectral bottleneck. Additionally, embedding exact conservative fluxes into the loss topology via finite-volume informed architectures could mitigate the unconstrained physical drift. Whether these architectural alternatives fully resolve the spectral barrier identified here remains an open question warranting systematic investigation. At a deeper geometric level, structure-preserving neural architectures — in particular, those designed to respect Hamiltonian or Lagrangian mechanics via symplectic integration layers \cite{Chen2020symplectic, Matsubara2021} — may offer an alternative route to resolving the spectral barrier by embedding the geometric structure of the PDE solution manifold directly into the architecture, rather than penalising its violation through soft constraints.

The Atlas presented here provides a reproducible, quantitative foundation for that investigation.
